%! Piecewise-Defined Functions — Unit 2, Lesson 2.6 — Guided Notes
\documentclass[10pt]{article}
\usepackage{atda-article}
\usepackage{atda-boxes}
\newcommand{\TallMath}[1]{$\displaystyle #1\rule[-1.4em]{0pt}{3.2em}$}
\setlength{\workrowsep}{6pt}

\begin{document}

\pageheader{Unit 2, Lesson 2.6}{Guided Notes}
% NAMESTRIP: no \namedateperiod here. The student writes their name once, on the
% cover this component is stapled behind. See references/conventions.md.

\vspace{0.15in}

\begin{objectivebox}
\small
By the end of this lesson, I will be able to\dots
\begin{enumerate}[leftmargin=1.5em, itemsep=2pt, topsep=3pt]
  \item read a graph that takes \textbf{more than one rule}, and evaluate it by first asking
        \emph{which piece owns this input} and then using that piece's rule and only that one;
  \item explain why a piecewise-defined function is still \textbf{one} function --- every input in
        the domain belongs to exactly one piece, so every input still has exactly one output;
  \item use the \textbf{open and closed circles} at a boundary to say what the function's value is
        there, and to write a range with a bracket or a parenthesis;
  \item decide whether a graph is \textbf{continuous} at a boundary, and tell a \emph{corner} from a
        \emph{break};
  \item find all of Unit 2's characteristics on a piecewise graph --- domain and range, zeros and
        intercepts, increasing/decreasing/constant intervals, absolute maximum and minimum, end
        behavior, and asymptotes.
\end{enumerate}
\end{objectivebox}

\vspace{0.10in}

\boxguard
\begin{vocabbox}
\small
Fill in each term as we build it together. Every one of them is read straight off a graph.
\termblanklong{Piecewise-defined function}
\termblanklong{Piece}
\termblanklong{Boundary point}
\termblanklong{Open circle / closed circle}
\termblanklong{Continuous / discontinuous}
\end{vocabbox}

\vspace{0.10in}

\boxguard[12]
\begin{hookbox}
\small
\textbf{Two students, one paycheck, and \$60 that nobody can find.} A garden center pays \$12 an hour
for the first $40$ hours of a week, and \$18 an hour for every hour beyond $40$. A student worked
$45$ hours. \emph{What is the pay for that week?}

\smallskip
Student A: \textbf{``$45$ hours at \$12 is \$540, and there were $5$ overtime hours at \$18, which is
another \$90. So \$630.''}

\smallskip
Student B: \textbf{``$40$ hours at \$12 is \$480, and $5$ hours at \$18 is \$90. So \$570.''}

\smallskip
\textbf{Here is the question: which hours does each rate apply to?} Both students used both rates and
neither made an arithmetic mistake. Decide who is right, and say exactly which hours the other one
paid twice.

\writelines{2}
\end{hookbox}

\vspace{0.10in}

\boxguard[30]
\begin{notesbox}{1. One Function, Written in Pieces}
\small
Last night's delivery fee took two sentences to describe: \emph{flat \$5 for the first three miles,
then \$2 a mile after that.} There is a name for a function like that, and there is a way to write it
down.

\vspace{4pt}
A \textbf{piecewise-defined function} is \blank{1.4cm} function described by more than one rule, with
each rule applying only on its own stretch of the \blank{2.2cm}. Each rule together with the stretch
of inputs it owns is called a \textbf{piece}.

\vspace{4pt}
\begin{center}
\begin{tikzpicture}
\begin{axis}[
    width=10.4cm, height=4.4cm,
    xlabel={$m$ (miles)}, ylabel={$F$ (fee, \$)},
    xmin=0, xmax=8, ymin=0, ymax=16,
    xtick={0,1,2,3,4,5,6,7,8}, ytick={0,5,10,15}, minor y tick num=4,
    grid=both, grid style={linegray!45}, minor grid style={linegray!30},
    axis lines=left,
    tick label style={font=\tiny}, label style={font=\scriptsize}]
  \addplot[royal, very thick] coordinates {(0,5) (3,5) (8,15)};
\end{axis}
\end{tikzpicture}
\end{center}

\vspace{-4pt}
{\scriptsize Last night's item 10: $F(m)$ is the fee in dollars for a delivery of $m$ miles, for
deliveries up to $8$ miles.}

\vspace{4pt}
Here is the same graph written in symbols. Read the brace as the word \emph{if}:
\[
  F(m) = \begin{cases}
    5 & \text{if } 0 \le m \le 3,\\[2pt]
    5 + 2(m-3) & \text{if } 3 < m \le 8.
  \end{cases}
\]

\vspace{2pt}
\textbf{Which piece owns which inputs?} Fill the right-hand column.

\vspace{2pt}
\renewcommand{\arraystretch}{1.4}
\begin{tabularx}{\linewidth}{@{}p{4.4cm} X@{}}
\rowcolor{mist}
\textbf{Rule} & \textbf{The inputs it owns} \\
\hline
$F(m)=5$ & \blank{6.0cm} \\
$F(m)=5+2(m-3)$ & \blank{6.0cm} \\
\end{tabularx}

\vspace{6pt}
\textbf{Evaluating takes one extra step, and it is always the same step.} Before you use a rule, ask
\textbf{which piece owns this input}. Then use that rule and \blank{2.4cm}.

\vspace{2pt}
\renewcommand{\arraystretch}{1.3}
\begin{tabular}{@{}|l|c|c|c|c|c|@{}}
\hline
$m$ (miles) & $1$ & $3$ & $4$ & $6$ & $8$ \\ \hline
Which piece? (1st or 2nd) & \blank{1.1cm} & \blank{1.1cm} & \blank{1.1cm} & \blank{1.1cm} &
    \blank{1.1cm} \\ \hline
$F(m)$ (\$) & \blank{1.1cm} & \blank{1.1cm} & \blank{1.1cm} & \blank{1.1cm} & \blank{1.1cm} \\
\hline
\end{tabular}

\vspace{6pt}
\textbf{The sentence that runs the whole lesson.} A piecewise function is \blank{2.2cm} function, not
several. The pieces are not allowed to overlap, so every input in the domain belongs to exactly
\blank{1.4cm} piece --- which means every input still has exactly \blank{2.4cm}. That is Lesson 2.1's
definition of a function, and it has not changed.
\end{notesbox}

\vspace{0.10in}

\boxguard[30]
\begin{notesbox}{2. Where Each Rule Lives --- and Why the Paycheck Was Wrong}
\small
Here is the garden center's pay function. $P(h)$ is one week's pay in dollars for $h$ hours worked,
and the week is capped at $50$ hours.

\begin{center}
\begin{tikzpicture}
\begin{axis}[
    width=10.4cm, height=4.8cm,
    xlabel={$h$ (hours worked in the week)}, ylabel={$P$ (pay, \$)},
    xmin=0, xmax=50, ymin=0, ymax=700,
    xtick={0,5,10,15,20,25,30,35,40,45,50}, ytick={0,120,240,360,480,600},
    minor y tick num=1, minor x tick num=4,
    grid=both, grid style={linegray!45}, minor grid style={linegray!30},
    axis lines=left,
    tick label style={font=\tiny}, label style={font=\scriptsize}]
  \addplot[royal, very thick] coordinates {(0,0) (40,480) (50,660)};
\end{axis}
\end{tikzpicture}
\end{center}

\vspace{-2pt}
\[
  P(h) = \begin{cases}
    12h & \text{if } 0 \le h \le 40,\\[2pt]
    480 + 18(h-40) & \text{if } 40 < h \le 50.
  \end{cases}
\]

\vspace{2pt}
\textbf{Read these five values off the graph.}

\vspace{2pt}
\renewcommand{\arraystretch}{1.3}
\begin{tabular}{@{}|l|c|c|c|c|c|@{}}
\hline
$h$ (hours) & $10$ & $20$ & $40$ & $45$ & $50$ \\ \hline
$P(h)$ (\$) & \blank{1.2cm} & \blank{1.2cm} & \blank{1.2cm} & \blank{1.2cm} & \blank{1.2cm} \\
\hline
\end{tabular}

\vspace{6pt}
\textbf{Now settle the hook.} Student A said \$630 for a $45$-hour week. The graph at $h=45$ says
\$\blank{1.2cm}. \textbf{Student \blank{0.8cm} answered the question.} Student A charged the \$12 rate
to \emph{all $45$} hours and then added overtime on top, so hours $41$ through $45$ got paid
\blank{2.0cm}. \textbf{A rule applies only on the stretch of inputs it \blank{2.0cm}} --- that is the
whole content of the word \emph{piecewise}, and it is the error to watch for all year.

\vspace{6pt}
\textbf{What happens at the boundary $h=40$?} Check both rules there: $12(40)=480$, and the second
rule starts at $480$ too. The two rules \blank{1.8cm} at the boundary, so the graph has no break ---
you can trace it from end to end without lifting your pencil. A graph like that is called
\blank{2.6cm}.

\vspace{4pt}
\textbf{And here is the distinction to write down.} At $h=40$ the graph turns a sharp
\blank{1.8cm} --- the line gets steeper because the pay rate went up. That is a change of
\emph{slope}, not a hole in the graph. \textbf{A corner is not a \blank{1.6cm}.}

\vspace{6pt}
\textbf{Every characteristic from this unit still works.} Fill the table --- one graph, all of Unit 2.

\vspace{2pt}
\renewcommand{\arraystretch}{1.4}
\begin{tabularx}{\linewidth}{@{}p{5.0cm} X@{}}
\rowcolor{mist}
\textbf{Characteristic} & \textbf{For $P$} \\
\hline
Domain (Lesson 2.1) & \blank{6.4cm} \\
Range (Lesson 2.1) & \blank{6.4cm} \\
Zeros and $y$-intercept (2.4) & \blank{6.4cm} \\
Increasing / decreasing /\newline constant (2.4) & \blank{6.4cm} \\
Absolute max and min (2.4) & \blank{6.4cm} \\
End behavior, asymptotes (2.5) & \blank{6.4cm} \\
\end{tabularx}
\end{notesbox}

\vspace{0.10in}

\boxguard[30]
\begin{notesbox}{3. When the Pieces Disagree --- Jumps, and What the Dots Mean}
\small
A shipping company charges by weight: \$6 for a package weighing up to and including $2$ pounds, \$10
for one over $2$ pounds up to and including $5$, and \$16 for one over $5$ pounds up to and including
$10$. Nothing in between --- the price changes all at once.

\begin{center}
\begin{tikzpicture}
\begin{axis}[
    width=10.4cm, height=4.6cm,
    xlabel={$w$ (pounds)}, ylabel={$S$ (shipping cost, \$)},
    xmin=0, xmax=10.6, ymin=0, ymax=20,
    xtick={0,1,2,3,4,5,6,7,8,9,10}, ytick={0,2,4,6,8,10,12,14,16,18},
    grid=both, grid style={linegray!45},
    axis lines=left,
    tick label style={font=\tiny}, label style={font=\scriptsize}]
  \addplot[royal, very thick] coordinates {(0,6) (2,6)};
  \addplot[royal, very thick] coordinates {(2,10) (5,10)};
  \addplot[royal, very thick] coordinates {(5,16) (10,16)};
  \addplot[royal, only marks, mark=*, mark size=2pt] coordinates {(2,6) (5,10) (10,16)};
  \addplot[royal, only marks, mark=*, mark size=2pt, mark options={fill=white}]
    coordinates {(0,6) (2,10) (5,16)};
\end{axis}
\end{tikzpicture}
\end{center}

\vspace{-4pt}
{\scriptsize A \textbf{filled} circle means that endpoint \emph{is} part of the graph. An
\textbf{open} circle means it is not.}

\vspace{4pt}
\textbf{The question the dots exist to answer: what is $S(2)$?} Look straight up from $w=2$. There
are two circles there --- one filled at \$6 and one open at \$10 --- so $S(2)=\$$\blank{1.2cm}. There
is exactly \blank{1.4cm} filled circle above every input, and there has to be, because a function
gives each input \blank{2.4cm}.

\vspace{4pt}
\textbf{Fill in each cost.}

\vspace{2pt}
\renewcommand{\arraystretch}{1.3}
\begin{tabular}{@{}|l|c|c|c|c|c|c|@{}}
\hline
$w$ (pounds) & $1$ & $2$ & $3$ & $5$ & $7$ & $10$ \\ \hline
$S(w)$ (\$) & \blank{1.0cm} & \blank{1.0cm} & \blank{1.0cm} & \blank{1.0cm} & \blank{1.0cm} &
    \blank{1.0cm} \\
\hline
\end{tabular}

\vspace{6pt}
\textbf{Three things this graph does that last box's graph did not.}
\begin{enumerate}[leftmargin=1.5em, itemsep=3pt, topsep=3pt]
  \item \textbf{It breaks.} To trace it you must lift your pencil at $w=$ \blank{1.0cm} and at $w=$
        \blank{1.0cm}. A graph with a break in it is \blank{2.8cm}. The break itself is called a
        \textbf{jump}: here the cost jumps by \$4 and then by \$6.
  \item \textbf{Its range is a list, not an interval.} The only outputs this function ever produces
        are \blank{4.0cm}. No interval notation can say that --- you write the set, exactly as you did
        with the hot dogs in Lesson 2.1.
  \item \textbf{It is never increasing.} On each of its three stretches the graph is
        \blank{2.2cm}, and a stretch is the only place ``increasing'' can be measured. The cost does
        go up --- but it goes up \emph{at} the jumps, and a jump happens over \blank{2.6cm} at all.
\end{enumerate}

\vspace{4pt}
\textbf{And the dots decide a bracket or a parenthesis.} The first piece owns the inputs
\blank{3.0cm}: open at $0$ because a package with no weight is not a package, closed at $2$ because a
$2$-pound package pays \$6. That is Lesson 2.1's rule, back again --- a bracket \blank{2.0cm} the
endpoint and a parenthesis leaves it out.
\end{notesbox}

\vspace{0.10in}

\boxguard[30]
\begin{practicebox}
\small
A delivery drone takes off, climbs, cruises at a fixed altitude, then descends to the ground. The
graph shows $A(t)$, its altitude in feet, $t$ seconds after take-off. It takes \emph{three} pieces.

\begin{center}
\begin{tikzpicture}
\begin{axis}[
    width=10.6cm, height=4.6cm,
    xlabel={$t$ (seconds after take-off)}, ylabel={$A$ (altitude, feet)},
    xmin=0, xmax=14, ymin=0, ymax=60,
    xtick={0,1,2,3,4,5,6,7,8,9,10,11,12,13,14}, ytick={0,10,20,30,40,50},
    minor y tick num=4,
    grid=both, grid style={linegray!45}, minor grid style={linegray!30},
    axis lines=left,
    tick label style={font=\tiny}, label style={font=\scriptsize}]
  \addplot[royal, very thick, domain=0:5, samples=60]{20*x-2*x^2};
  \addplot[royal, very thick] coordinates {(5,50) (9,50) (14,0)};
\end{axis}
\end{tikzpicture}
\end{center}

\begin{enumerate}[leftmargin=1.6em, itemsep=6pt, topsep=2pt]
  \item Read these altitudes off the graph, and name which piece --- climb, cruise, or descent ---
        owns each input.
\smallskip

\renewcommand{\arraystretch}{1.3}
\begin{tabular}{@{}|l|c|c|c|c|c|c|@{}}
\hline
$t$ (seconds) & $0$ & $2$ & $5$ & $8$ & $11$ & $14$ \\ \hline
$A(t)$ (feet) & \blank{1.0cm} & \blank{1.0cm} & \blank{1.0cm} & \blank{1.0cm} & \blank{1.0cm} &
    \blank{1.0cm} \\ \hline
Which piece? & \blank{1.0cm} & \blank{1.0cm} & \blank{1.0cm} & \blank{1.0cm} & \blank{1.0cm} &
    \blank{1.0cm} \\
\hline
\end{tabular}

  \item Give the domain and the range of $A$, then give its zeros and its $y$-intercept. Say what
        each zero means about the flight.

\writelines{3}

  \item Name the interval on which $A$ is increasing, the interval on which it is constant, and the
        interval on which it is decreasing. Then give the absolute maximum --- \emph{value and
        location}. Read the location carefully.

\writelines{3}

  \item Is this graph continuous? Answer yes or no with a reason. Then answer the Lesson 2.5
        questions: what is the end behavior of $A$, and does it have any asymptotes?

\writelines{3}
\end{enumerate}
\end{practicebox}

\end{document}
